\documentclass[11pt, a4paper]{article}

% Gemeinsame Style-Definitionen (TEXINPUTS muss templates/ enthalten — siehe scripts/build_tex.py)
% bfw-style.tex — gemeinsame Style-Definitionen für alle Handouts/Aufgabenhefte
% Voraussetzung: Kompilation mit xelatex (wegen fontspec)

\usepackage[ngerman]{babel}
% Babel-ngerman macht " zu einem aktiven Shorthand (z.B. für "a -> ä). In unseren
% Texten verwenden wir " als normales Zeichen — daher Shorthand komplett ausschalten.
\shorthandoff{"}
\usepackage{geometry}
\geometry{a4paper, top=2.4cm, bottom=2.4cm, left=2.3cm, right=2.3cm,
          headheight=15pt, headsep=0.7cm, footskip=1.1cm}

\usepackage{fontspec}
% Fallback-Kette: Source Sans Pro -> Calibri -> Segoe UI -> Latin Modern Sans
% (Latin Modern ist immer mit MiKTeX vorhanden)
\IfFontExistsTF{Source Sans Pro}{%
  \setmainfont{Source Sans Pro}[Ligatures=TeX]%
}{\IfFontExistsTF{Calibri}{%
  \setmainfont{Calibri}[Ligatures=TeX]%
}{\IfFontExistsTF{Segoe UI}{%
  \setmainfont{Segoe UI}[Ligatures=TeX]%
}{%
  \setmainfont{Latin Modern Sans}[Ligatures=TeX]%
}}}
\IfFontExistsTF{Source Code Pro}{%
  \setmonofont{Source Code Pro}[Scale=0.92]%
}{\IfFontExistsTF{Consolas}{%
  \setmonofont{Consolas}[Scale=0.92]%
}{%
  \setmonofont{Latin Modern Mono}[Scale=0.92]%
}}

% Gesperrte Versalien für Labels/Titelbalken (Kapitälchen-Ersatz, funktioniert
% mit jeder OpenType-Schrift — Latin Modern Sans hat keine echten Kapitälchen)
\newcommand{\bfwlabelfont}{\footnotesize\bfseries\addfontfeatures{LetterSpace=8.0}}

\usepackage{xcolor}
% Farbpalette: hell, freundlich, modern
\definecolor{bfwPrimary}{HTML}{0EA5A4}    % Petrol/Türkis - Primär
\definecolor{bfwPrimaryDark}{HTML}{0F766E} % dunkler Petrol für Headings
\definecolor{bfwAccent}{HTML}{F59E0B}     % Amber - Akzent
\definecolor{bfwSuccess}{HTML}{10B981}    % Grün - Erfolg / Lösung
\definecolor{bfwWarn}{HTML}{F97316}       % Orange - Achtung
\definecolor{bfwInk}{HTML}{0F172A}        % fast Schwarz
\definecolor{bfwInkSoft}{HTML}{475569}    % gedämpftes Grau
\definecolor{bfwBgSoft}{HTML}{F8FAFC}     % off-white
\definecolor{bfwBgTeal}{HTML}{ECFEFF}     % helles Türkis
\definecolor{bfwBgAmber}{HTML}{FEF3C7}    % helles Gelb
\definecolor{bfwBgRose}{HTML}{FFE4E6}     % helles Rosé für Eselsbrücken
\definecolor{bfwTabHead}{HTML}{E4F3F2}    % zarte Kopfzeilen-Hinterlegung für Tabellen

\usepackage[colorlinks=true, linkcolor=bfwPrimaryDark, urlcolor=bfwPrimaryDark]{hyperref}
\usepackage{lastpage}
\usepackage{enumitem}
\setlist[itemize]{leftmargin=*, itemsep=2pt, topsep=2pt}
\setlist[enumerate]{leftmargin=*, itemsep=2pt, topsep=2pt}

\usepackage[most]{tcolorbox}
\usepackage{booktabs}
\usepackage{colortbl}
\usepackage{tikz}
\usetikzlibrary{decorations.pathmorphing, positioning, shapes.geometric, arrows.meta}

% ---------------------------------------------------------------------------
% Einheitliches Tabellen-Design
%   - etwas mehr Zeilenluft, dezente graue Linien (wirkt auch mit \hline ruhiger)
%   - Kopfzeile einer Tabelle bei Bedarf zart hinterlegen: \rowcolor{bfwTabHead}
% ---------------------------------------------------------------------------
\renewcommand{\arraystretch}{1.25}
\arrayrulecolor{bfwInkSoft!50}
\setlength{\heavyrulewidth}{1pt}
\setlength{\lightrulewidth}{0.5pt}

% ---------------------------------------------------------------------------
% Boxen — klare farbige Titelbalken mit gesperrten Versal-Labels
% (Titel bleibt per Option überschreibbar: \begin{merksatz}[title={...}])
% ---------------------------------------------------------------------------
\tcbset{bfwbox/.style={%
  enhanced, breakable,
  boxrule=0.5pt, arc=2.5pt,
  left=10pt, right=10pt, top=8pt, bottom=8pt,
  toptitle=4pt, bottomtitle=4pt,
  titlerule=0pt,
  fonttitle=\bfwlabelfont,
  coltitle=white
}}

% Merkkästchen — wichtige Definition / Kernpunkt
\newtcolorbox{merksatz}[1][]{%
  bfwbox,
  colback=bfwBgTeal,
  colframe=bfwPrimaryDark,
  colbacktitle=bfwPrimaryDark,
  title={MERKE},
  #1
}

% Eselsbrücke — Mnemoniks
\newtcolorbox{eselsbruecke}[1][]{%
  bfwbox,
  colback=bfwBgRose,
  colframe=bfwAccent!85!black,
  colbacktitle=bfwAccent!85!black,
  title={ESELSBRÜCKE},
  #1
}

% Beispiel-Box
\newtcolorbox{beispiel}[1][]{%
  bfwbox,
  colback=bfwBgSoft,
  colframe=bfwInkSoft,
  colbacktitle=bfwInkSoft,
  title={BEISPIEL},
  #1
}

% Lösungs-Box (für Aufgabenhefte)
\newtcolorbox{loesung}[1][]{%
  bfwbox,
  colback=bfwSuccess!7,
  colframe=bfwSuccess!65!black,
  colbacktitle=bfwSuccess!65!black,
  title={LÖSUNG},
  #1
}

% Achtung-Box — typische Fehler / Stolperfallen
\newtcolorbox{achtung}[1][]{%
  bfwbox,
  colback=bfwWarn!6,
  colframe=bfwWarn!85!black,
  colbacktitle=bfwWarn!85!black,
  title={ACHTUNG},
  #1
}

% Formel-Box — für zentrale Formeln (groß, ruhig, ohne Titelbalken)
\newtcolorbox{formelbox}[1][]{%
  enhanced,
  colback=bfwBgTeal,
  colframe=bfwPrimaryDark,
  boxrule=1pt, arc=3pt,
  left=10pt, right=10pt, top=6pt, bottom=6pt,
  #1
}

% Glossar-Eintrag
\newcommand{\glossareintrag}[2]{%
  \par\noindent\textbf{\color{bfwPrimaryDark}#1} \enspace #2\par\smallskip
}

% ---------------------------------------------------------------------------
% Abschnitts-Design: farbiges Nummern-Quadrat + feine Linie unter \section,
% dezent farbige \subsection/\subsubsection
% ---------------------------------------------------------------------------
\usepackage{titlesec}
\newcommand{\bfwSecBadge}[1]{%
  \begin{tikzpicture}[baseline={([yshift=-0.62em]n.center)}]
    \node[fill=bfwPrimaryDark, text=white, font=\normalsize\bfseries,
          minimum width=1.55em, minimum height=1.55em, inner sep=1pt,
          rounded corners=1.5pt] (n) {#1};
  \end{tikzpicture}}
\titleformat{\section}
  {\Large\bfseries\color{bfwPrimaryDark}}
  {\bfwSecBadge{\thesection}}{0.6em}{}
  [\vspace{-0.2em}{\color{bfwPrimary!60}\titlerule[0.8pt]}]
\titleformat{\subsection}
  {\large\bfseries\color{bfwPrimaryDark}}
  {\textcolor{bfwPrimary}{\thesubsection}}{0.5em}{}
\titleformat{\subsubsection}
  {\normalsize\bfseries\color{bfwInk}}
  {\textcolor{bfwPrimary}{\thesubsubsection}}{0.4em}{}
\titlespacing*{\section}{0pt}{1.6em}{1em}
\titlespacing*{\subsection}{0pt}{1.2em}{0.5em}

% TikZ-Stil "skizzenhaft" — etwas weicher als Rough.js, aber stabil
\tikzset{
  roughbox/.style = {
    draw=bfwPrimaryDark, thick, fill=bfwBgTeal,
    rounded corners=4pt, inner sep=6pt
  },
  roughaccent/.style = {
    draw=bfwAccent, thick, fill=bfwBgAmber,
    rounded corners=4pt, inner sep=6pt
  },
  roughline/.style = {
    draw=bfwInkSoft, thick, -{Stealth[length=2.5mm]}
  }
}

% Bullet-Symbole (bewusst kein FontAwesome — vermeidet Font-Installations-Probleme)
\newcommand{\faLightbulb}{$\bullet$}
\newcommand{\faBookmark}{$\bullet$}

% ---------------------------------------------------------------------------
% Kopf-/Fußzeile
%   Kopf links:  Wortlogo „Wirtschaftsfachwirt"
%   Kopf rechts: Dokument-Kurztext, gesetzt via \kopfzeile{...}
%   Fuß links:   © Can Siebert · 2026 — Fuß rechts: Seite X von Y
%   Titelseite (Seite 1) bleibt ohne Kopf-/Fußzeile (\handouttitel setzt
%   \thispagestyle{empty}).
% ---------------------------------------------------------------------------
\usepackage{fancyhdr}
\newcommand{\bfwKopfzeileText}{}
\newcommand{\kopfzeile}[1]{\renewcommand{\bfwKopfzeileText}{#1}}
\pagestyle{fancy}
\fancyhf{}
\fancyhead[L]{\small\bfseries\color{bfwPrimaryDark}Wirtschaftsfachwirt}
\fancyhead[R]{\small\color{bfwInkSoft}\bfwKopfzeileText}
\renewcommand{\headrulewidth}{0.6pt}
\renewcommand{\headrule}{{\color{bfwPrimary}\hrule width\headwidth height 0.6pt}}
\fancyfoot[L]{\small\color{bfwInkSoft}\copyright\ Can Siebert\,\textperiodcentered\,2026}
\fancyfoot[R]{\small\color{bfwInkSoft}Seite \thepage\ von \pageref*{LastPage}}
% Auch \thispagestyle{plain} (z.B. durch \maketitle) soll leer bleiben:
\fancypagestyle{plain}{\fancyhf{}\renewcommand{\headrulewidth}{0pt}\renewcommand{\headrule}{}}

% ---------------------------------------------------------------------------
% \handouttitel{Obertitel}{Untertitel}{Kurs-Label}{Termin-Zeile}
%   Einheitlicher professioneller Titelblock: Dachzeile, farbige Headline,
%   Untertitel, farbige Trennlinie und dezente Meta-Zeile (Kurs/Termin/Dozent).
%   Ersetzt \title/\maketitle. Direkt nach \begin{document} aufrufen.
% ---------------------------------------------------------------------------
\newcommand{\handouttitel}[4]{%
  \thispagestyle{empty}%
  \begingroup
  \setlength{\parindent}{0pt}%
  {\bfwlabelfont\color{bfwPrimary}WIRTSCHAFTSFACHWIRT (IHK)\par}
  \vspace{0.55em}
  {\huge\bfseries\color{bfwPrimaryDark}#1\par}
  \vspace{0.5em}
  {\large\color{bfwInkSoft}#2\par}
  \vspace{0.9em}
  {\color{bfwPrimary}\rule{\linewidth}{2pt}}\par
  \vspace{0.55em}
  {\small\color{bfwInkSoft}#3\ \,\textperiodcentered\,\ #4\ \,\textperiodcentered\,\ Dozent: Can Siebert\par}
  \endgroup
  \vspace{1.4em}%
}


\kopfzeile{Handout BWL Lektion 4 \textperiodcentered{} Teil 2}

\begin{document}
\handouttitel{Lektion~4 \textperiodcentered{} Teil~2: Konzentration \& Fusionskontrolle}%
  {Konzern, Fusion und Trust --- und wo das Kartellrecht Grenzen zieht}%
  {1.4 Unternehmenszusammenschlüsse}%
  {Sa, 27.02.2027 \textperiodcentered{} Session 2}

\begin{merksatz}[title={\bfseries Worum geht es in Teil 2?}]
In Teil~1 haben Sie die Kooperationsformen kennengelernt, bei denen die Partner
selbstständig bleiben. Diese zweite Session behandelt die \textbf{Konzentration}:
Beim \textbf{Konzern} geht die wirtschaftliche Selbstständigkeit verloren, bei
\textbf{Fusion} und \textbf{Trust} zusätzlich die rechtliche. Außerdem ordnen Sie
Zusammenschlüsse nach ihrer \textbf{Richtung} (horizontal, vertikal, diagonal) ein,
vergleichen die Ziele von Kooperation und Konzentration und lernen die
\textbf{Fusionskontrolle} durch Bundeskartellamt und EU-Kommission als Grenze der
Konzentration kennen. Die Session schließt mit dem Kursrückblick (Rahmenplan 1.1--1.4)
und Prüfungshinweisen --- es ist der letzte Kurstermin.
\end{merksatz}

\begin{beispiel}[title={Hinweis zur Literatur}]
Der Textband-Auszug zu diesem Thema wird nachgereicht. Bis dahin dienen die beiden
Teil-Handouts als Skriptersatz; ergänzend: Übungsband Betriebswirtschaft,
Übungen 54--58 (mit Lösungen).
\end{beispiel}

\tableofcontents
\vspace{1em}
\hrule
\vspace{1em}

\section{Anknüpfung an Teil 1: die Kernabgrenzung}

Zur Erinnerung aus Session 1: \textbf{Kooperation} = freiwillige Zusammen\emph{arbeit}
rechtlich und wirtschaftlich selbstständiger Unternehmen (ARGE/Konsortium,
Interessengemeinschaft, Joint Venture, strategische Allianz, Verband --- verboten nur
das Kartell, \S\,1 GWB). \textbf{Konzentration} = Zusammen\emph{fassung} von Unternehmen
unter Verlust der wirtschaftlichen Selbstständigkeit --- beim Konzern bleibt die
rechtliche Selbstständigkeit bestehen, bei Fusion und Trust geht auch sie verloren.
Genau diese Konzentrationsformen sind jetzt unser Thema.

\section{Formen der Konzentration}

\subsection{Konzern}

Ein \textbf{Konzern} ist die Zusammenfassung eines \emph{herrschenden} und eines oder
mehrerer \emph{abhängiger} Unternehmen unter \textbf{einheitlicher Leitung}
(\S\,18 AktG). Die Konzernunternehmen bleiben \textbf{rechtlich selbstständig}, verlieren
aber ihre \textbf{wirtschaftliche Selbstständigkeit}.

\begin{itemize}
  \item \textbf{Muttergesellschaft (Obergesellschaft):} übt die einheitliche Leitung aus
    --- meist über eine \emph{Mehrheitsbeteiligung} an den Töchtern; eine reine
    \emph{Holding} beschränkt sich auf das Halten und Steuern der Beteiligungen.
  \item \textbf{Tochtergesellschaften:} rechtlich selbstständige Gesellschaften
    (typisch GmbH oder AG), die den Weisungen der Mutter folgen.
  \item \textbf{Beherrschungsvertrag} (\S\,291 AktG): Unternehmensvertrag, mit dem sich
    eine Gesellschaft der Leitung eines anderen Unternehmens unterstellt --- die Mutter
    erhält ein direktes \emph{Weisungsrecht} gegenüber dem Vorstand der Tochter; häufig
    kombiniert mit einem Gewinnabführungsvertrag.
  \item \textbf{Unterordnungskonzern} (Regelfall: Mutter--Tochter) vs.
    \textbf{Gleichordnungskonzern} (Unternehmen unter einheitlicher Leitung, ohne dass
    eines vom anderen abhängig ist).
\end{itemize}

\subsection{Fusion (Verschmelzung)}

Die \textbf{Fusion} ist die engste Form des Zusammenschlusses: Mindestens ein beteiligtes
Unternehmen gibt seine \textbf{rechtliche und wirtschaftliche Selbstständigkeit} auf
(geregelt im Umwandlungsgesetz). Zwei Wege:

\begin{itemize}
  \item \textbf{Verschmelzung durch Aufnahme:} Unternehmen B überträgt sein Vermögen auf
    Unternehmen A; B erlischt (A + B $\rightarrow$ A).
  \item \textbf{Verschmelzung durch Neugründung:} A und B übertragen ihr Vermögen auf
    eine neu gegründete Gesellschaft C; A und B erlöschen (A + B $\rightarrow$ C).
\end{itemize}

\subsection{Trust}

Historischer, v.\,a. aus den USA stammender Begriff für die \textbf{vollständige
Verschmelzung} ehemals selbstständiger Unternehmen zu einer wirtschaftlichen und
rechtlichen Einheit mit marktbeherrschender Stellung. Die beteiligten Unternehmen
verlieren rechtliche \emph{und} wirtschaftliche Selbstständigkeit; Ziel ist die
Monopolisierung ganzer Märkte --- daher der Begriff „Antitrust-Recht`` für das
Wettbewerbsrecht.

\begin{eselsbruecke}
Intensitäts-Treppe der Konzentration: \textbf{K}onzern (rechtlich noch selbstständig)
$\rightarrow$ \textbf{F}usion (ein Unternehmen erlischt) $\rightarrow$ \textbf{T}rust
(alles verschmilzt zur beherrschenden Einheit): \textbf{K}ommt \textbf{f}ast
\textbf{t}otal zusammen.
\end{eselsbruecke}

\section{Richtungen des Zusammenschlusses}

Unabhängig von der Form (Kooperation oder Konzentration) unterscheidet man, \emph{welche}
Unternehmen sich verbinden:

\begin{itemize}
  \item \textbf{Horizontal:} Unternehmen der \emph{gleichen} Branche und
    \emph{gleichen} Produktions- oder Handelsstufe --- z.\,B. zwei Brauereien, zwei
    Großküchenhersteller. Ziele: Marktanteil, Größenvorteile, Rationalisierung.
  \item \textbf{Vertikal:} Unternehmen \emph{aufeinanderfolgender} Wirtschaftsstufen ---
    \emph{vorgelagert} (Richtung Rohstoff/Zulieferer, „Rückwärtsintegration``: Brauerei +
    Mälzerei) oder \emph{nachgelagert} (Richtung Absatz/Kunde, „Vorwärtsintegration``:
    Brauerei + Getränkegroßhandel). Ziele: Sicherung von Beschaffung und Absatz.
  \item \textbf{Diagonal} (auch \emph{anorganisch}, \emph{lateral},
    \emph{konglomerat}): \emph{branchenfremde} Unternehmen ohne
    Leistungszusammenhang --- z.\,B. Maschinenbauer + Versicherung. Ziele:
    Risikostreuung (Diversifikation), Kapitalanlage, neue Geschäftsfelder.
\end{itemize}

\begin{beispiel}[title={Beispiel aus dem Übungsband (Übung 56): Großküchenhersteller}]
Ein Großküchenhersteller verbindet sich \emph{vertikal vorgelagert} mit seinem
Spülmaschinen-Zulieferer, \emph{vertikal nachgelagert} mit einem Kantinenbetreiber
(Abnehmer), \emph{horizontal} mit einem anderen Großküchenhersteller und
\emph{diagonal} mit einer Geschäftsbank (branchenfremd).
\end{beispiel}

\section{Ziele im Vergleich und Grenzen: die Fusionskontrolle}

\subsection{Ziele der Kooperation vs. Ziele der Konzentration}

\begin{center}
\begin{tabular}{p{6.8cm} p{6.8cm}}
\toprule
\textbf{Kooperation --- „gemeinsam stärker, aber frei``} & \textbf{Konzentration --- „Einheit und Kontrolle``}\\
\midrule
Synergien in Teilbereichen nutzen (F\&E, Einkauf, Vertrieb) & umfassende Rationalisierung unter einheitlicher Leitung\\
Großprojekte gemeinsam stemmen (ARGE, Konsortium) & Marktmacht durch Marktanteilsaddition\\
Risiko einzelner Projekte teilen & Risikostreuung durch Diversifikation (konglomerat)\\
Märkte gemeinsam erschließen (Joint Venture, Allianz) & Sicherung von Beschaffung und Absatz durch Integration\\
Selbstständigkeit und Flexibilität bleiben erhalten & volle Kontrolle, aber Verlust der Eigenständigkeit\\
\bottomrule
\end{tabular}
\end{center}

\subsection{Fusionskontrolle durch Bundeskartellamt und EU-Kommission}

Konzentration kann den Wettbewerb dauerhaft ausschalten --- deshalb unterliegen größere
Zusammenschlüsse der \textbf{präventiven Fusionskontrolle}:

\begin{itemize}
  \item \textbf{Deutschland} (\S\S\,35\,ff. GWB): Zusammenschlüsse müssen beim
    \textbf{Bundeskartellamt} (Bonn) \emph{vor} dem Vollzug \emph{angemeldet} werden,
    wenn die beteiligten Unternehmen bestimmte Umsatzschwellen überschreiten (weltweit
    zusammen mehr als 500~Mio.~€; zudem im Inland ein Beteiligter mehr als 50~Mio.~€
    und ein weiterer mehr als 17,5~Mio.~€ Umsatz).
  \item \textbf{Untersagung:} Der Zusammenschluss wird verboten, wenn er wirksamen
    Wettbewerb \emph{erheblich behindern} würde --- insbesondere wenn eine
    \textbf{marktbeherrschende Stellung} entsteht oder verstärkt wird (Vermutung ab
    40\,\% Marktanteil eines Unternehmens). Möglich sind auch Freigaben unter
    \emph{Auflagen} (z.\,B. Verkauf von Unternehmensteilen).
  \item \textbf{Ministererlaubnis} (\S\,42 GWB): Ausnahmsweise kann das
    Bundeswirtschaftsministerium einen untersagten Zusammenschluss aus
    Gemeinwohlgründen dennoch erlauben.
  \item \textbf{EU-Ebene:} Bei \emph{gemeinschaftsweiter Bedeutung} (sehr hohe,
    europaweite Umsätze) prüft die \textbf{EU-Kommission} nach der
    EU-Fusionskontrollverordnung (FKVO) --- dann ist sie anstelle des Bundeskartellamts
    zuständig („One-Stop-Shop``).
\end{itemize}

\begin{merksatz}
Kartell\textbf{verbot} (Teil~1) richtet sich gegen \emph{Absprachen} (Verhalten),
Fusions\textbf{kontrolle} gegen \emph{Strukturen} (dauerhafte Konzentration).
\end{merksatz}

\section{Formeln \& Rechenbeispiele}

Zu Teil~2 gehören zwei Rechenwege (der dritte --- der Synergieeffekt --- steht im
Handout zu Teil~1). Prägen Sie sich jeweils Formel, Rechenweg und die prüfungsrelevante
Schlussfolgerung ein.

\subsection{Marktanteil (Bezug: Fusionskontrolle)}

\[
\textrm{Marktanteil} \;=\; \frac{\textrm{eigener Umsatz}}{\textrm{Gesamtumsatz des Marktes}} \times 100\,\%
\]

\begin{beispiel}[title={Rechenbeispiel: Marktanteile bei einer geplanten Fusion}]
Der Gesamtmarkt hat ein Volumen von 800~Mio.~€. Unternehmen U setzt 240~Mio.~€ um und
will Unternehmen W (120~Mio.~€ Umsatz) übernehmen.
\begin{enumerate}
  \item \textbf{Marktanteil U:} $240 \div 800 \times 100\,\% = 30\,\%$
  \item \textbf{Marktanteil W:} $120 \div 800 \times 100\,\% = 15\,\%$
  \item \textbf{Gemeinsamer Marktanteil nach der Fusion:} $30\,\% + 15\,\% = 45\,\%$
\end{enumerate}
\textbf{Schlussfolgerung:} $45\,\% > 40\,\%$ --- die
\emph{Marktbeherrschungsvermutung} des GWB greift; das Bundeskartellamt prüft den
Zusammenschluss besonders kritisch.
\end{beispiel}

\subsection{Beteiligungsquote im Konzern}

\[
\textrm{Beteiligungsquote} \;=\; \frac{\textrm{gehaltene Stimmrechte}}{\textrm{gesamte Stimmrechte}} \times 100\,\%
\]

\begin{beispiel}[title={Rechenbeispiel: Mehrheitsbeteiligung}]
Die M-AG hält 120.000 von insgesamt 200.000 Stimmrechten der T-GmbH.
\begin{enumerate}
  \item \textbf{Quote berechnen:} $120.000 \div 200.000 = 0{,}60$
  \item \textbf{In Prozent:} $0{,}60 \times 100\,\% = 60\,\%$
  \item \textbf{Bewerten:} $60\,\% > 50\,\%$ $\rightarrow$ \emph{Mehrheitsbeteiligung}
\end{enumerate}
\textbf{Schlussfolgerung:} Die M-AG kann Beschlüsse der T-GmbH durchsetzen ---
Abhängigkeit und Konzern werden vermutet; die T-GmbH verliert ihre wirtschaftliche
(nicht die rechtliche) Selbstständigkeit.
\end{beispiel}

\begin{merksatz}
Faustregeln zur Stimmrechtsquote: über 50\,\% = \emph{Mehrheitsbeteiligung}
(Beherrschung möglich), ab 75\,\% = qualifizierte Mehrheit (z.\,B. für
Satzungsänderungen), 100\,\% = hundertprozentige Tochter.
\end{merksatz}

\section{Kursabschluss: Rückblick und Prüfungsvorbereitung}

Mit dieser Session ist der Rahmenplan-Abschnitt 1.4 --- und damit der Kurs ---
abgeschlossen. Der Kursrückblick umfasst sieben Lektionen: VWL~1 (Grundbegriffe \&
Preisbildung), VWL~2 (Wettbewerb, Staatseingriffe \& VGR), VWL~3 (Konjunktur,
Wirtschaftspolitik \& Außenwirtschaft), BWL~1 und BWL~2 (betriebliche Funktionen),
BWL~3 (Existenzgründung \& Rechtsformen) und BWL~4 (Unternehmenszusammenschlüsse).

\begin{merksatz}[title={\bfseries Wichtig für die Prüfungsvorbereitung}]
Sie sollten sicher können: die Abgrenzung \emph{Kooperation vs. Konzentration} über die
rechtliche und wirtschaftliche Selbstständigkeit (Teil~1), die Definition des
\emph{Konzerns} (rechtlich selbstständig, einheitliche Leitung, Beherrschungsvertrag),
die \emph{Fusionswege} (Aufnahme/Neugründung), die drei \emph{Richtungen} mit
Praxisbeispielen, \emph{Kartellverbot} (\S\,1 GWB, Teil~1) und \emph{Fusionskontrolle}
(Bundeskartellamt/EU-Kommission) sowie die drei \emph{Rechenwege} (Marktanteil,
Beteiligungsquote, Synergieeffekt). Üben Sie mit dem Aufgabenheft, den beiden Quizzen
und den Übungen 54--58 des Übungsbands --- und nutzen Sie das Prüfungstraining BWL auf
dem Dashboard!
\end{merksatz}

\vspace{1em}
\hrule
\vspace{0.8em}
\noindent\textsf{\small Quellen: Rahmenplan Wirtschaftsbezogene Qualifikationen (DIHK), Abschnitt 1.4;
Übungsband Betriebswirtschaft (DIHK-Bildungs-GmbH), Übungen 54--58; GWB, AktG, UmwG (Grundzüge).
Textband-Auszug wird nachgereicht. Dies ist der letzte Kurstermin --- viel Erfolg bei der IHK-Prüfung!}

\end{document}
